\documentclass[11pt,a4paper]{article}

\usepackage[utf8]{inputenc}
\usepackage[italian]{babel}
\usepackage{amsmath}
\usepackage{amsfonts}
\usepackage{pgfplots}
\pgfplotsset{compat=1.18}
\usepackage{geometry}
\geometry{a4paper, margin=2cm}

\title{
    \textbf{Teoria dei Segnali (Modulo: Segnali Determinati e Aleatori)
    : Tesina Numero 1
}}

\author{Andrea D'Amico \\ \small{Ingegneria Informatica (L-8) | Università degli studi di Catania}}
\date{}

\begin{document}

\maketitle

\section{Introduzione e Modello Analitico}

L'obiettivo di questa relazione è l'analisi in frequenza e la successiva sintesi del segnale periodico $y(t)$ assegnato. Al fine di semplificare il calcolo dei coefficienti dello sviluppo in serie di Fourier, risulta conveniente sfruttare la linearità dell'operatore trasformata e scomporre il segnale in componenti elementari note.

Osservando l'andamento del segnale in un singolo periodo $T_0$, e centrando opportunamente l'asse dei tempi, è possibile modellare $y(t)$ come la sovrapposizione lineare di $y_1(t)$ (segnale $\Delta$) e $y_2(t)$ (segnale $\Pi$).

Questo è il modo più efficiente di calcolare le caratteristiche del segnale, in quanto le proprietà delle funzioni $\Delta$ e $\Pi$ sono già note.

I parametri calcolati a partire dal numero di lettere del cognome ("Damico" $\rightarrow N=6$) sono:
\begin{itemize}
    \item $A = \lceil \frac{N}{2} \rceil = \lceil \frac{6}{2} \rceil = 3 \text{ V}$
    \item $T_0 = \frac{1}{N} \text{ ms} = \frac{1}{6} \text{ ms}$
\end{itemize}

L'espressione analitica nel periodo $t \in \left[-\frac{T_0}{2}, \frac{T_0}{2}\right]$ risulta essere:
\begin{equation}
    y(t) = y_1(t) + y_2(t) = A \cdot \Delta\left(\frac{t}{T_0/2}\right) + \frac{A}{2} \Pi\left(\frac{t}{T_0/2}\right)
\end{equation}

\subsection{Sfruttamento delle Simmetrie}
Una proprietà fondamentale di questa scomposizione è che sia $\Delta(\cdot)$ che $\Pi(\cdot)$ sono \textbf{funzioni pari}. Di conseguenza, la loro somma $y(t)$ è una funzione pari. Questa simmetria garantisce, per il teorema di Fourier, che la serie sarà composta unicamente da termini cosinusoidali (parte immaginaria nulla).

\newpage

\subsection{Rappresentazione Grafica delle Componenti}
Di seguito viene illustrata la scomposizione grafica del segnale (con valori normalizzati $A=3, T_0=4$ per chiarezza visiva), mostrando prima le singole componenti e infine il segnale risultante.

\vspace{0.5cm}

\begin{center}
\begin{minipage}{0.48\textwidth}
\centering
\begin{tikzpicture}
\begin{axis}[
    width=1.1\textwidth,
    height=4cm,
    axis lines=middle,
    xmin=-3, xmax=3,
    ymin=-0.5, ymax=3.5,
    xlabel={$t$},
    ylabel={$y_1(t)$},
    grid=both,
    grid style={dashed, gray!30},
    thick,
    title={Componente $\Delta$}
]
\addplot[
    domain=-2:2, 
    samples=100, 
    color=blue
]
{3 * max(0, 1 - abs(x)/2)};
\end{axis}
\end{tikzpicture}
\end{minipage}
\hfill
\begin{minipage}{0.48\textwidth}
\centering
\begin{tikzpicture}
\begin{axis}[
    width=1.1\textwidth,
    height=4cm,
    axis lines=middle,
    xmin=-3, xmax=3,
    ymin=-0.5, ymax=2.0,
    xlabel={$t$},
    ylabel={$y_2(t)$},
    grid=both,
    grid style={dashed, gray!30},
    thick,
    title={Componente $\Pi$}
]
\addplot[
    domain=-2:2, 
    samples=200, 
    color=red
]
{(abs(x) <= 1) ? 1.5 : 0};
\end{axis}
\end{tikzpicture}
\end{minipage}
\end{center}

\vspace{0.8cm}

\begin{center}
\begin{tikzpicture}
\begin{axis}[
    width=0.9\textwidth,
    height=9cm,
    axis lines=middle,
    xmin=-3, xmax=3,
    ymin=-0.5, ymax=5.0,
    xlabel={$t$},
    ylabel={Ampiezza},
    legend pos=outer north east,
    legend cell align={left},
    grid=both,
    grid style={dashed, gray!30},
    thick,
    title={Somma delle componenti $y(t) = y_1(t) + y_2(t)$}
]

% Componente 1: \Delta (y1)
\addplot[
    domain=-2:2, 
    samples=100, 
    color=blue, 
    dashed
]
{3 * max(0, 1 - abs(x)/2)};
\addlegendentry{$y_1(t)$}

% Componente 2: \Pi (y2)
\addplot[
    domain=-2:2, 
    samples=200, 
    color=red, 
    dashed
]
{(abs(x) <= 1) ? 1.5 : 0};
\addlegendentry{$y_2(t)$}

% Segnale Originale (y)
\addplot[
    domain=-2:2, 
    samples=200, 
    color=violet,
    ultra thick
]
{(abs(x) <= 1 ? 1.5 : 0) + 3 * max(0, 1 - abs(x)/2)};
\addlegendentry{$y(t) = y_1(t) + y_2(t)$}

\end{axis}
\end{tikzpicture}
\end{center}

Come si evince dai grafici, la somma tra il fronte di discesa di $\Pi$ e la pendenza costante di $\Delta$ in $t = \pm \frac{T_0}{4}$ genera esattamente la discontinuità di ampiezza $\frac{A}{2}$ richiesta dalle specifiche del segnale originale.

\vspace{0.8cm}
\begin{center}
\begin{tikzpicture}
\begin{axis}[
    width=0.9\textwidth,
    height=7cm,
    axis lines=middle,
    xmin=-0.12, xmax=0.12,
    ymin=-0.5, ymax=5.0,
    xlabel={$t$ (ms)},
    ylabel={Ampiezza (V)},
    xtick={-0.08333, -0.04166, 0, 0.04166, 0.08333},
    xticklabels={$-\frac{1}{12}$, $-\frac{1}{24}$, $0$, $\frac{1}{24}$, $\frac{1}{12}$},
    scaled x ticks=false,
    tick label style={font=\footnotesize},
    grid=both,
    grid style={dashed, gray!30},
    thick,
    title={Segnale $y(t)$ con Parametri Reali ($A=3\text{ V}, T_0=1/6\text{ ms}$)}
]

% T_0/2 = 1/12 ms. Per ricondurci alla scala normalizzata, valutiamo le funzioni su x*12
\addplot[
    domain=-0.1:0.1, 
    samples=300, 
    color=violet,
    ultra thick
]
{(abs(x*12) <= 1 ? 1.5 : 0) + 3 * max(0, 1 - abs(x*12))};

\end{axis}
\end{tikzpicture}
\end{center}

\newpage
\section{Sviluppo in Serie di Fourier}
% 1. Calcolare lo sviluppo in serie di Fourier (ampiezza e fase)

\newpage
\section{Spettro delle Ampiezze e delle Fasi}
% 2. Disegnare lo spettro dele ampiezze e delle fasi delle prime 10 armoniche

\newpage
\section{Ricostruzione e Simulazione Software}
% 2. Ricostruire attraverso GnuRadio (o altro ambiente sw) il segnale come somma delle prime 10 armoniche del segnale
\subsection{Sintesi con 10 Armoniche e Verifica nel Tempo}
% 3. Verificare che il segnare nel tempo restituito da GnuRadio e’ uguale a quello assegnato

\subsection{Verifica dello Spettro}
% 4. Verificare che lo spettro delle ampiezze restituto da GnuRadio sia uguale a quello calcolato in precedenza

\subsection{Sintesi con 3 Armoniche ed Errore di Troncamento}
% 5. Ricostruire attraverso GnuRadio il segnale come somma delle prime 3 armoniche e commentare l'errore rispetto a quello con 10 armoniche

\newpage
\section{Calcolo della Potenza e Banda del Segnale}
% 7. Calcolare quante armoniche sono necessarie per ricostruire un segnale che abbia una potenza media maggiore del 99,95% di quella del segnale di partenza

\newpage
\section{Analisi del Segnale Sfasato}
% 8. Riscrivere l’espressione del segnale sfasato in figura 1 e rifare i punti da 1 a 4

\end{document}